% ═══════════════════════════════════════════════════════════════════════
%  preambule.tex — Réglages, couleurs et raccourcis communs du diaporama
%  du laboratoire « Introduction à LaTeX ».
%  Ce fichier est autonome : il ne dépend d'aucun fichier du livre.
% ═══════════════════════════════════════════════════════════════════════

% ── Encodage, langue, polices (lualatex, comme le livre) ───────────────
\usepackage{fontspec}
\usepackage[french]{babel}

\setmainfont{TeX Gyre Termes}
\IfFontExistsTF{Fira Sans}
  {\setsansfont{Fira Sans}[Scale=MatchLowercase]}
  {\setsansfont{TeX Gyre Heros}[Scale=MatchLowercase]}
\usefonttheme{professionalfonts}
\renewcommand{\familydefault}{\sfdefault}

\usepackage{amsmath}   % avant unicode-math (sinon \underbrace est cassé)
\usepackage{mathtools}
\usepackage{unicode-math}
\IfFontExistsTF{NewCMMath-Book.otf}
  {\setmathfont{NewCMMath-Book}}
  {\setmathfont{Latin Modern Math}}

\usepackage{booktabs}
\usepackage{array}
\usepackage{graphicx}
\usepackage{adjustbox}
\usepackage{siunitx}
\sisetup{output-decimal-marker = {,}, exponent-product = \times,
         group-separator = {\,}, group-minimum-digits = 4}
\usepackage[skins,breakable]{tcolorbox}
\usepackage{tikz}
\usetikzlibrary{calc, arrows.meta, angles, quotes, patterns, positioning,
                decorations.markings, backgrounds, fit, shapes.geometric}

% ── Palette (identique à celle du livre, livrejs.cls) ──────────────────
\definecolor{sectioncolor}{RGB}{28, 63, 110}   % bleu foncé — titres
\definecolor{thmboxbg}{RGB}{228, 238, 252}     % bleu pâle — énoncés
\definecolor{coveraccent}{RGB}{196, 142, 58}   % or chaud — accent
\definecolor{secboxbg}{RGB}{253, 240, 219}     % ocre pâle — « Pour la classe »
\definecolor{steelblue4}{HTML}{36648B}         % figures : objet de départ
\definecolor{firebrick3}{HTML}{CD3333}         % figures : objet image
\definecolor{demobg}{RGB}{247, 247, 247}       % fond neutre — « au tableau »
\definecolor{bgpage}{RGB}{253, 253, 252}       % fond de diapositive, légèrement chaud

% ── Thème beamer : sobre, cohérent avec le livre ────────────────────────
\usecolortheme{default}
\setbeamercolor{palette primary}{bg=sectioncolor, fg=white}
\setbeamercolor{palette secondary}{bg=sectioncolor, fg=white}
\setbeamercolor{palette tertiary}{bg=sectioncolor, fg=white}
\setbeamercolor{title}{fg=sectioncolor}
\setbeamercolor{frametitle}{bg=sectioncolor, fg=white}
\setbeamercolor{structure}{fg=sectioncolor}
\setbeamercolor{section in toc}{fg=sectioncolor}
\setbeamercolor{subsection in toc}{fg=black}
\setbeamercolor{background canvas}{bg=bgpage}
\setbeamercolor{footline}{fg=gray, bg=bgpage}
\setbeamercolor{block title}{bg=sectioncolor, fg=white}
\setbeamercolor{block body}{bg=thmboxbg, fg=black}
\setbeamercolor{item}{fg=sectioncolor}
\setbeamercolor{itemize item}{fg=coveraccent}
\setbeamercolor{itemize subitem}{fg=sectioncolor}
\setbeamertemplate{itemize item}{\raise0.3ex\hbox{\small$\blacktriangleright$}}
\setbeamertemplate{itemize subitem}{\raise0.3ex\hbox{\footnotesize$\bullet$}}
\setbeamertemplate{navigation symbols}{}
\setbeamerfont{frametitle}{series=\bfseries}
\setbeamertemplate{frametitle}{%
  \vskip-1ex
  \begin{beamercolorbox}[wd=\paperwidth,ht=3.8ex,dp=1.3ex,leftskip=1em,rightskip=1em]{frametitle}
    \bfseries\strut\insertframetitle
  \end{beamercolorbox}
  \vskip3pt
}
\setbeamertemplate{footline}{%
  \leavevmode
  \hbox{\begin{beamercolorbox}[wd=\paperwidth,ht=2.4ex,dp=1ex,leftskip=1em,rightskip=1em]{footline}
    {\scriptsize\color{gray}Écrire des mathématiques avec \LaTeX
    \hfill \insertshorttitle{} \hfill \insertframenumber/\inserttotalframenumber}
  \end{beamercolorbox}}%
}
\setbeamertemplate{blocks}[rounded][shadow=false]
\setbeamertemplate{section in toc}{%
  \inserttocsectionnumber.~\inserttocsection\par\vskip0.15em}

% ── Page de section : diapositive de transition ────────────────────────
\AtBeginSection[]{%
  \begin{frame}[plain,c]
    \begin{center}
      {\color{coveraccent}\rule{0.25\textwidth}{1.2pt}}\\[10pt]
      {\usebeamerfont{frametitle}\usebeamercolor[fg]{title}%
       \Large\bfseries\color{sectioncolor}\insertsectionhead}\\[10pt]
      {\color{sectioncolor}\rule{0.45\textwidth}{0.5pt}}
    \end{center}
  \end{frame}
}
\AtBeginSubsection[]{%
  \begin{frame}[plain,c]
    \begin{center}
      {\large\color{gray}\insertsectionhead}\\[4pt]
      {\Large\bfseries\color{coveraccent}\insertsubsectionhead}
    \end{center}
  \end{frame}
}

% ── Raccourcis mathématiques repris du livre (autonomes) ────────────────
\newcommand{\Plan}{\Pi}
\newcommand{\Transf}{\mathcal{T}}
\newcommand{\Aff}{\mathcal{A}}
\newcommand{\Sim}{\mathcal{S}}
\newcommand{\Iso}{\mathcal{I}}
\newcommand{\mat}[4]{\begin{pmatrix}#1 & #2\\ #3 & #4\end{pmatrix}}
\newcommand{\col}[2]{\begin{pmatrix}#1\\ #2\end{pmatrix}}
\newcommand{\coliii}[3]{\begin{pmatrix}#1\\ #2\\ #3\end{pmatrix}}
\newcommand{\norme}[1]{\lVert #1 \rVert}
\newcommand{\Aire}{\operatorname{Aire}}
\newcommand{\dete}{\operatorname{det}}
\renewcommand{\vec}[1]{\mathbf{#1}}
\DeclareMathOperator{\R}{\mathbb{R}}
\DeclareMathOperator{\Rs}{\mathbb{R}^{\star}}
\newcommand{\intff}[2]{\mathopen{[}{#1},{#2}\mathclose{]}}
\newcommand{\intoo}[2]{\mathopen{]}{#1},{#2}\mathclose{[}}
\newcommand{\defi}[1]{{\color{coveraccent}#1}}
\newcommand{\intfo}[2]{\mathopen{[}{#1},{#2}\mathclose{[}}
\DeclareMathOperator{\N}{\mathbb{N}}
\DeclareMathOperator{\Ns}{\mathbb{N}^{\star}}
\DeclareMathOperator{\Q}{\mathbb{Q}}

% ── Raccourcis propres au chapitre 2 (identiques à macros-livre.tex) ────
\newcommand{\dE}{d_E}                         % distance euclidienne
\newcommand{\dT}{d_T}                         % taxi-distance (Manhattan)
\newcommand{\dHam}{d_{\mathrm{Ham}}}          % distance de Hamming
\newcommand{\dH}{d_H}                         % distance de Hausdorff
\newcommand{\boule}[2]{B(#1, #2)}             % boule ouverte
\newcommand{\Vois}[2]{#1^{\oplus #2}}         % epsilon-voisinage (Minkowski)
\newcommand{\Comp}[1]{\mathcal{K}(#1)}        % compacts non vides de E
\newcommand{\comp}[1]{\mathopen{}#1\mathclose{}^{\mathsf{c}}}

% Images bitmap (portraits, figures Python) reprises du livre.
\graphicspath{{exemples/}}

% ── Styles TikZ communs à toutes les figures (identiques au livre) ──────
\tikzset{
  axe/.style        = {-{Stealth[length=2mm]}, gray!70, thin},
  quadrillage/.style= {gray!22, very thin},
  fig/.style        = {fill=steelblue4!22, draw=steelblue4, thick,
                       line join=round},
  figim/.style      = {fill=firebrick3!22, draw=firebrick3, thick,
                       line join=round},
  figaux/.style     = {fill=coveraccent!22, draw=coveraccent!85!black,
                       thick, line join=round},
  pt/.style         = {circle, fill=sectioncolor, inner sep=1.1pt},
  ptim/.style       = {circle, fill=firebrick3, inner sep=1.1pt},
  fleche/.style     = {-{Stealth[length=2.4mm]}, sectioncolor, thick},
  etiq/.style       = {font=\scriptsize, inner sep=1.5pt},
  droite/.style     = {sectioncolor!70, thick, dashed},
}
\newcommand{\repere}[4]{%
  \draw[quadrillage] (#1,#3) grid (#2,#4);
  \draw[axe] (#1,0) -- (#2+0.35,0) node[etiq, below] {$x$};
  \draw[axe] (0,#3) -- (0,#4+0.35) node[etiq, left] {$y$};}
\newcommand{\motifF}{(0,0) -- (0.3,0) -- (0.3,0.55) -- (0.6,0.55) --
  (0.6,0.75) -- (0.3,0.75) -- (0.3,0.85) -- (0.75,0.85) -- (0.75,1) --
  (0,1) -- cycle}

% ── Boîtes d'énoncés (même esprit que le livre : fond bleu pâle) ────────
\tcbset{
  boxenonce/.style={
    enhanced, breakable, colback=thmboxbg, colframe=sectioncolor,
    boxrule=0.6pt, arc=2pt, left=9pt, right=9pt, top=6pt, bottom=6pt,
    fonttitle=\bfseries\sffamily, coltitle=white,
    colbacktitle=sectioncolor, attach boxed title to top left,
    boxed title style={arc=1pt, boxrule=0pt, left=6pt, right=6pt,
      top=3pt, bottom=3pt},
    title style={sharp corners},
  },
  boxclasse/.style={
    enhanced, breakable, colback=secboxbg, colframe=coveraccent,
    boxrule=0.6pt, arc=2pt, left=9pt, right=9pt, top=6pt, bottom=6pt,
    fonttitle=\bfseries\sffamily, coltitle=white,
    colbacktitle=coveraccent, attach boxed title to top left,
    boxed title style={arc=1pt, boxrule=0pt, left=6pt, right=6pt,
      top=3pt, bottom=3pt},
  },
  boxrem/.style={
    enhanced, breakable, colback=white, colframe=steelblue4,
    boxrule=0pt, leftrule=3pt, arc=0pt, left=9pt, right=9pt,
    top=5pt, bottom=5pt,
    colbacktitle=white, coltitle=steelblue4, toptitle=2pt,
    bottomtitle=0pt, fonttitle=\bfseries\sffamily,
  },
  boxoutil/.style={
    enhanced, breakable, colback=white, colframe=steelblue4,
    boxrule=0.6pt, arc=2pt, left=9pt, right=9pt, top=6pt, bottom=6pt,
    fonttitle=\bfseries\sffamily, coltitle=white,
    colbacktitle=steelblue4, attach boxed title to top left,
    boxed title style={arc=1pt, boxrule=0pt, left=6pt, right=6pt,
      top=3pt, bottom=3pt},
  },
  boxtableau/.style={
    enhanced, breakable, colback=demobg, colframe=firebrick3,
    boxrule=0pt, leftrule=4pt, arc=0pt, left=10pt, right=10pt,
    top=7pt, bottom=7pt,
  },
}

% Environnements : titre = « Type N.M — description », fourni en argument.
\newtcolorbox{Enonce}[1]{boxenonce, title={#1}}
\newtcolorbox{Classe}[1]{boxclasse, title={Pour la classe --- #1}}
\newtcolorbox{Outil}[1]{boxoutil, title={Outil --- #1}}
\newtcolorbox{Rem}[1]{boxrem, title={#1}}

% « Démonstration au tableau » : marqueur + rappel très court de l'idée.
\newenvironment{AuTableau}[1][Démonstration au tableau]{%
  \begin{tcolorbox}[boxtableau]
  \bfseries\sffamily\color{firebrick3}\MakeUppercase{#1}\par\nointerlineskip
  \vspace{4pt}\normalfont\rmfamily\upshape\mdseries\color{black}
}{%
  \end{tcolorbox}
}

% Petit renvoi discret vers le numéro du livre, en fin de bloc de titre.
\newcommand{\refnotes}[1]{{\scriptsize\color{gray} (notes de cours, #1)}}

% Insère une figure (tikzpicture ou image) en la contraignant à la fois
% en largeur et en hauteur, sans distorsion : \FigBox[hauteur max]{contenu}.

% ═══════════════════════════════════════════════════════════════════════
%  Environnements propres au laboratoire LaTeX
% ═══════════════════════════════════════════════════════════════════════
\tcbuselibrary{listings}
\definecolor{codebg}{RGB}{246, 248, 251}
\definecolor{vertpiece}{RGB}{58, 125, 68}
\definecolor{promptcol}{RGB}{92, 74, 138}      % violet sobre : les prompts
\definecolor{promptbg}{RGB}{245, 242, 250}

% Code LaTeX : commandes en bleu, commentaires en gris, jamais de gras.
\lstdefinestyle{latexcode}{
  language=[LaTeX]TeX, basicstyle=\ttfamily\scriptsize,
  keywordstyle=\color{sectioncolor}, texcsstyle=*\color{sectioncolor},
  commentstyle=\color{gray}, columns=fullflexible, keepspaces=true,
  upquote=true, showstringspaces=false, breaklines=true,
  moretexcs={documentclass, usepackage, title, author, date, maketitle,
    section, subsection, item, frac, dfrac, sqrt, num, qty, sisetup,
    question, part, parts, choice, CorrectChoice, printanswers, gradetable,
    numpoints, numquestions, pointpoints, hqword, hpword, hsword, htword,
    firstpageheader, runningheader, firstpagefooter, runningfooter,
    fillwithlines, makeemptybox, stretch, vspace, includegraphics, caption,
    centering, text, mathbb, cdot, times, leq, geq, neq, approx, pm, pi,
    infty, sin, cos, log, left, right, vec, quad, textsuperscript, hline,
    draw, node, solutiontitle, renewcommand, newcommand, label, ref, lhead,
    rhead, cfoot, fancyhf, pagestyle, thepage, numpages, textwidth,
    centi, metre, squared, minute, alpha, beta, theta, Delta, in, to,
    displaystyle, ldots, thispagestyle, tikz},
  literate={é}{{é}}1 {è}{{è}}1 {ê}{{ê}}1 {à}{{à}}1 {â}{{â}}1
           {ç}{{ç}}1 {ù}{{ù}}1 {û}{{û}}1 {ô}{{ô}}1 {î}{{î}}1 {ï}{{ï}}1
           {É}{{É}}1 {À}{{À}}1 {«}{{«}}1 {»}{{»}}1 {œ}{{œ}}1 {’}{{’}}1
           {—}{{—}}1 {–}{{–}}1
}

\tcbset{
  codestyle/.style={
    enhanced, colback=codebg, colframe=sectioncolor!35, boxrule=0.4pt,
    arc=2pt, left=6pt, right=6pt, top=1pt, bottom=1pt,
    listing options={style=latexcode},
    fonttitle=\sffamily\footnotesize, coltitle=white,
    colbacktitle=sectioncolor!80, attach boxed title to top left={yshift=-2mm, xshift=3mm},
    boxed title style={arc=1pt, boxrule=0pt, left=4pt, right=4pt, top=1pt, bottom=1pt},
  },
}
% Code seul : \begin{CodeTeX}[title=...] ... \end{CodeTeX}
\newtcblisting{CodeTeX}[1][]{codestyle, listing only, #1}
% Code et rendu côte à côte (le code est réellement compilé dans la diapo).
\newtcblisting{CodeRendu}[1][]{codestyle, listing side text,
  lefthand ratio=0.56, sidebyside gap=4mm, sidebyside align=center seam,
  segmentation style={draw=sectioncolor!35, dashed},
  fontlower=\footnotesize, #1}

% Prompt type pour l'IA : bulle de conversation sobre.
\newcommand{\iconebulle}{%
  \tikz[baseline=-0.7ex]{\fill[white, rounded corners=1pt] (0,-0.08) rectangle (0.34,0.16);
    \fill[white] (0.06,-0.07) -- (0.02,-0.17) -- (0.14,-0.07) -- cycle;}}
\newtcolorbox{Prompt}[1]{
  enhanced, colback=promptbg, colframe=promptcol, boxrule=0pt,
  leftrule=3pt, arc=0pt, outer arc=0pt, left=9pt, right=9pt, top=5pt,
  bottom=5pt, fontupper=\footnotesize\sffamily\raggedright,
  fonttitle=\sffamily\footnotesize, coltitle=white, colbacktitle=promptcol,
  title={\iconebulle\hspace{0.45em}Prompt type --- #1},
  attach boxed title to top left={yshift=-2mm, xshift=4mm},
  boxed title style={arc=1pt, boxrule=0pt, left=5pt, right=5pt, top=2pt, bottom=2pt},
  before upper={\setlength{\parskip}{3pt}},
}
% Un champ à compléter dans un prompt : [niveau], [sujet], etc.
\newcommand{\champ}[1]{{\color{promptcol}[#1]}}
% Une commande LaTeX citée dans le texte : \cmd{section} → \section
\newcommand{\cmd}[1]{\texttt{\textbackslash #1}}

% Exercice en laboratoire : même esprit que « Pour la classe ».
\newtcolorbox{Exercice}[1]{boxclasse, title={À vous --- #1}}

\newcommand{\FigBox}[2][0.56\textheight]{%
  \begin{center}
    \adjustbox{max width=\linewidth, max height=#1}{#2}
  \end{center}}
